%%%%%%%%%%%%%%%%%%%%%%%%
%%% 1-genesis/01.tex %%%
%%%%%%%%%%%%%%%%%%%%%%%%

\Chapter{1}

\Verse{1}
{
    \hP{בְּרֵאשִׁית בָּרָא אֱלֹהִים אֵת הַשָּׁמַיִם}
    \hP{וְאֵת הָאָרֶץ}
}
{
    \eP{In the beginning of {\God}’s creating the skies
        and the earth}
}
{
    \aB{
        Ok so, according to the \textsc{DRWFEROTMSOHK}%
        \footnote{Abbreviation for \textsc{The Dynamic Read-Writable Free Encyclopedic Repository of the Modern State of Human Knowledge}, hereafter abbreviated as ``the \textsc{Lord},'' to save space.}

        \begin{center}
        The word \hBP{אלהים} (Elhim) is a ``defective spelling'' of the word \hBP{אלוהים} (Elohim)
        \end{center}

        And if you look up that word, it says

        \begin{center}
        The word \hBP{אלוהים} (Elohim) is from the plural of \hBP{אלוה} (Eloh) which means ``god.''
        \end{center}

        And if you look around for some clarification, you'll find sentences like

        \begin{center}
        This word is sometimes grammatically plural and sometimes singular.
        \end{center}

        And you may also stumble upon a sentence mentioning that particularly devout believers

        \begin{center}
        ...may prefer to write and say \hBP{אלוקים} (Elokim), so as not to abuse the name of God.
        \end{center}

        So we're off to a pretty weird start.
    }
}

\Verse{2}
{
    \hP{וְהָאָרֶץ הָיְתָה תֹהוּ וָבֹהוּ}
    \hP{וְחֹשֶׁךְ עַל פְּנֵי תְהוֹם וְרוּחַ אֱלֹהִים מְרַחֶפֶת עַל פְּנֵי}
    \hP{הַמָּיִם}
}
{
    \eP{when the earth had been shapeless
        and formless, and darkness was on the face of the deep,
        and {\God}’s spirit was hovering on the face of the water}
}
{
    \aB{
        Ok so apparently we didn't know how to read the past perfect in biblical Hebrew until the 1970s.

        So it turns out the earth was already around when this Elohim character started creating it.

        Hence the ``had been'' in this line, and the odd grammar in the last.

        Just make a note of that and let's keep moving.
    }
}

\Verse{3}
{
    \hP{וַיֹּאמֶר אֱלֹהִים יְהִי אוֹר וַיְהִי אוֹר}
}
{
    \eP{{\GodInit} said, “Let there be light.” And there was light.}
}
{

    \aB{
        Words create stuff. Like this book.

        Not a huge surprise. At least to us authors.%
        \footnote{The bible, it turns out, was written by an author. Or rather, several.}

        I mean obviously.
    }
}

\Verse{4}
{
    \hP{וַיַּרְא אֱלֹהִים אֶת הָאוֹר כִּי טוֹב}
    \hP{וַיַּבְדֵּל אֱלֹהִים בֵּין הָאוֹר}
    \hP{וּבֵין הַחֹשֶׁךְ}
}
{
    \eP{And {\God} saw the light, that it was good,
        and {\God} separated between the light
        and the darkness.}
}
{
    \aB{
        Ok so after he creates light, he decides he did a good job.

        Then after that, he gets the darkness out of it.

        That's sort of like ``cleaning the light'' I guess.
    }
}

\Verse{5}
{
    \hP{וַיִּקְרָא אֱלֹהִים לָאוֹר יוֹם}
    \hP{וְלַחֹשֶׁךְ קָרָא לָיְלָה וַיְהִי עֶרֶב}
    \hP{וַיְהִי בֹקֶר יוֹם אֶחָד}
}
{
    \eP{And {\God} called the light “day” and called the darkness “night.”
        And there was evening, and there was morning: one day.}
}
{
    \aB{
        Ok so he names the light, and the dark light that he cleaned out of the original light.

        He calls them day and night. Now remember, words create stuff.

        So immediately he's affected by this decision and it's suddenly day two.

        It's like if you ``invent murder'' and then suddenly there's a dead body.

        Because obviously. Hilarious.
    }
}

\Verse{6}
{
    \hP{וַיֹּאמֶר אֱלֹהִים יְהִי רָקִיעַ בְּתוֹךְ הַמָּיִם}
    \hP{וִיהִי מַבְדִּיל בֵּין מַיִם לָמָיִם}
}
{
    \eP{And {\God} said, “Let there be a space within the water,
        and let it separate between water
        and water.”}
}
{
    \aB{
        It's not super clear what's going on here yet.

        But he says some words.
    }
}

\Verse{7}
{
    \hP{וַיַּעַשׂ אֱלֹהִים אֶת הָרָקִיעַ}
    \hP{וַיַּבְדֵּל בֵּין הַמַּיִם אֲשֶׁר מִתַּחַת לָרָקִיעַ}
    \hP{וּבֵין הַמַּיִם אֲשֶׁר מֵעַל לָרָקִיעַ}
    \hP{וַיְהִי כֵן}
}
{
    \eP{And {\God} made the space, and it separated between the water that was
        under the space and the water that was above the space.
        And it was so.}
}
{
    \aB{
        Ok so it seems like the words didn't create the stuff this time.

        Because here after saying ``Let there be,'' he still has to make the thing.

        So he takes some of the water and pushes it up above the other water.

        The ``up water'' is the sky.

        Now wait, that's not as dumb as it sounds.

        I mean (1) it's blue up there, and (2) water comes down.

        So we must live in a middle bit between the ``up water'' and the ``down water.''

        Completely reasonable so far.
    }
}

\Verse{8}
{
    \hP{וַיִּקְרָא אֱלֹהִים לָרָקִיעַ שָׁמָיִם}
    \hP{וַיְהִי עֶרֶב וַיְהִי בֹקֶר יוֹם שֵׁנִי}
}
{
    \eP{And {\God} called the space “skies.”
        And there was evening, and there was morning: a second day.}
}
{
    \aB{
        The ``up water'' isn't the sky.

        The sky is the bit in the middle, between the ``up water'' and ``down water.''

        No land yet. He's not done.

        Also the day ends before he got a chance to say ``it was good.''

        Hopefully he'll say it twice tomorrow and get back on schedule.
    }
}

\Verse{9}
{
    \hP{וַיֹּאמֶר אֱלֹהִים יִקָּווּ הַמַּיִם מִתַּחַת הַשָּׁמַיִם אֶל מָקוֹם}
    \hP{אֶחָד וְתֵרָאֶה הַיַּבָּשָׁה וַיְהִי כֵן}
}
{
    \eP{And {\God} said, “Let the waters be concentrated under the skies into one
        place, and let the land appear.”\\
        And it was so.}
}
{
    \aB{
        So he moves some of the down water horizontally so it all piles up in one place.

        Then land appears. He didn't make the land, it was already there, like we saw in line two.
    }
}

\Verse{10}
{
    \hP{וַיִּקְרָא אֱלֹהִים לַיַּבָּשָׁה אֶרֶץ}
    \hP{וּלְמִקְוֵה הַמַּיִם קָרָא יַמִּים}
    \hP{וַיַּרְא אֱלֹהִים כִּי טוֹב}
}
{
    \eP{And {\God} called the land “earth,”
        and called the concentration of the waters “seas.”
        And {\God} saw that it was good.}
}
{
    \aB{
        Now remember, we already met the ``earth'' in line one.

        But now he calls the pile of water over there ``seas.''

        And the land over here is ``earth.''

        Same Hebrew word as the first line.

        It's spelled \hBP{ארץ} (arts).

        That's earth.

        Also he said it's good again.
    }
}

\Verse{11}
{
    \hP{וַיֹּאמֶר אֱלֹהִים תַּדְשֵׁא הָאָרֶץ דֶּשֶׁא עֵשֶׂב מַזְרִיעַ זֶרַע}
    \hP{עֵץ פְּרִי עֹשֶׂה פְּרִי לְמִינוֹ אֲשֶׁר זַרְעוֹ בוֹ עַל הָאָרֶץ}
    \hP{וַיְהִי כֵן}
}
{
    \eP{And {\God} said, “Let the earth generate plants, vegetation that produces
        seed, fruit trees, each making fruit of its own kind, which has its seed
        in it, on the earth. And it was so:”}
}
{
    \aB{At this point we're starting to notice that the author of this section
        is pretty verbose.}
}

\Verse{12}
{
    \hP{וַתּוֹצֵא הָאָרֶץ דֶּשֶׁא עֵשֶׂב מַזְרִיעַ זֶרַע לְמִינֵהוּ}
    \hP{וְעֵץ עֹשֶׂה פְּרִי אֲשֶׁר זַרְעוֹ בוֹ לְמִינֵהוּ}
    \hP{וַיַּרְא אֱלֹהִים כִּי טוֹב}
}
{
    \eP{The earth brought out plants, vegetation that produces seeds of its own
        kind, and trees that make fruit that each has seeds of its own kind in
        it. And {\God} saw that it was good.}
}
{
    \aB{
        Really verbose. That'll be relevant later.

        Also, that's the second ``it was good'' today.
    }
}

\Verse{13}
{
    \hP{וַיְהִי עֶרֶב וַיְהִי בֹקֶר יוֹם שְׁלִישִׁי}
}
{
    \eP{And there was evening, and there was morning: a third day.}
}
{
    \aB{So we're all caught up from missing yesterday.}
}

\Verse{14}
{
    \hP{וַיֹּאמֶר אֱלֹהִים יְהִי מְאֹרֹת בִּרְקִיעַ הַשָּׁמַיִם לְהַבְדִּיל}
    \hP{בֵּין הַיּוֹם וּבֵין הַלָּיְלָה וְהָיוּ לְאֹתֹת}
    \hP{וּלְמוֹעֲדִים וּלְיָמִים וְשָׁנִים}
}
{
    \eP{And {\God} said, “Let there be lights in the space of the skies to
        distinguish between the day and the night,
        and they will be for signs and for appointed times
        and for days and years.}
}
{
    \aB{
                Ok so, at this point he decides to make two big sky balls.

                Apparently the light and dark-light have just been sort of flying around on their own until now,
                causing day and night all by themselves, and that's confusing.

                So he makes two sky balls.

                To distinguish day and night. I guess that was hard before now.

                The sky balls are for ``appointed times.''

                Like holidays and stuff. Or like dinner time and bed time.

                They're also for days and years by the way.

                Every time the day ball goes around, that's a day.

                Every time the night ball goes around, that's a month.

                You may notice the girls bleed about once per moon too.

                So it seems like the moon and girls are related somehow.%
                \footnote{
                \aR{\footnotesize Editor's Note:}
                \aC{\footnotesize
                    The original text contains here a lengthy footnote,
                    which we found to be both 1. here, and 2. concerning,
                    concerning Chinese mythology, and beginning
                }
                Cf., Chinese mythology,
                \aC{\footnotesize
                    and ending much later on a rather undignified note, a note which,
                    to put it mildly, is largely unrelated to Chinese mythology,
                    and which \textsc{We} shall not reproduce here, in the interest
                    of maintaining the integrity of this anthology
                    }%
                \aB{\footnotesize
				where the universe develops out of a primordial chaos energy organized
				into the cycles of Yin and Yang, and where the former is unsurprisingly
				associated with the quiet, hidden, subtle, coy, lunar feminine,
				while the latter is associated with the more Loud Bright
				Wham-Bam-Pow-Zap-Zam type of Not-Nearly-So-Subtle-Hushed-or-Hidden guy
				s\redacted{hit}\aR{\footnotesize tuff}
				like wearing a big foam finger and a hat with two beer can holders
				on either side with plastic tube type straws going down so you can 
				drink your beer while you cheer for your favorite Let's Go Sports Team™ while
				painted head to toe in the Official Colors® of Said Sports Team©®™
				and after a few too many drinks in a bout of celebratory elevated in-group
				revelatory enthusiasm you whip your \redacted{dick} out and do that helicopter dance
				that guys sometimes do where they spin \redacted{their dick around}
				like the propeller of the aforementioned helicopulatory vehicle
				for maybe three or four seconds max before it's captured by one
				of the many audience-view-cameras just when the guy in the control room
				(he and his lawyers insist, in no uncertain terms, by accident)
				switched the focus of the big stadium view screen to precisely that camera,
				and in so doing, thereby broadcast, for a few shining moments (which
				the Chairperson of the National \textsc{Whatever-Sport-This-Is} Board Of
				Trustees and Franchisees describes as ``the most disgraceful three or
				four seconds in National \textsc{W/e} STI League history'') the aforementioned
				super-fan's member to over ten thousand fans and members not to mention the
				members of the Fran Qizi Memorial Franchisees' and Trustees'
				MVFaT Box Club seated up in the box seats with the aristocratic
				old-money-type investors and the new-money Chinese family into
				which they've recently married, who decided at that point that
				they'd seen quite enough and it was therefore time for them to
				leave in a huff.
				That stuff.
				That's Yang.
				Real man \redacted{shit} my guy.
				Yang's like that.
				And Yin's the other way.
				And you need both to \aR{\footnotesize be}
				f\redacted{uck}\aR{\footnotesize ruitful} and make life
				so remember, seek the harmony of the dual.
				BDSM: Binary Duality, Sun and Moon. There's balance in all things.
				And sure this bit down here's a bit long and hard to take in but hey,
				every long hard Yang needs a short thin easy Yin at the end of the day.
				Not easy for everyone. But easy for Yang. She's Yin and he's in.
				And on that note I'm out, so don't forget, this hard lengthy part
				down here you're looking at now is Yang.
				So Yin's like the moon, and Yang's like all this,
				and Yin likes all this no matter what she might say
				(and even though (as you know) this long Yang part of the
				corpus down here's a bit big and rough and raw compared to
				the soft smoother sort of Yin
				but\redacted{t} but hey at the end of the day---%
				when it's the appointed time to get the moon out and get a little looney,
				whether or not you're in the mood for all this hard-to-take-in length
				we've got down here, we'll make sure you take it all in even if we have
				to wrestle till sunrise cuz when one or two come or comes to either end of the day---%
				you gotta love a good duel because god\redacted{dammit}
				\aR{\footnotesize is good, and duality} it's fun\aR{\footnotesize damental})
				and also naturally, sort of by process of elimination, the sun.%
					{\!\!\dn{0.333ex}{
					\chineseB{
					    \up{0.666ex}{\size{9}{\cK{月}}}%
					    \size{9}{\cR{曰}}%
					    \dn{0.666ex}{\size{9}{\cB{日}}}%
					}}}
				}}

                That's why we name them things like \textsc{Luna}.

                Every time the night ball goes around twelve times, that's a year.

                Every time the girls bleed twelve times, that's a year too.

                So twelve is a bible number.

                Something is ``bible'' if it's universal.

                Like if it's about all the humans everywhere.

                For example, if someone wrote something super mysterious about ``the water of life''

                And if they said it was ``clear as crystal''%
                \footnote{%
                    Translator's Note: Crystals are more of a cloudy white,
                    and often far from clear.
                }%

                And there was a whole river of the stuff flowing down out of the middle of a great ``street''%
                \footnote{
                \label{fn:street}%
                    Translator's Note: πλατείας in the original, meaning:
                    1. wide; broad; wide spread; spread out (in the adjectival form) or
                    2. wide open space; dish, cup, concave vessel-like space;
                       or sometimes, apparently, street (in the nounal.)
                    Synonyms include ευρύς (meaning wide) and φαρδύς (meaning wide, broad.)
                }%

                Down from a great throne%
                \footnote{
                \label{fn:seat}%
                    Translator's Note: θρόνου in the original, meaning both ``throne'' and ``seat'',
                    but translated here as the former in order to ensure that the meaning of this passage
                    is clear as crystal.
                }
                through the center of the spread wide apart area\refootnote{fn:street}

                And from the center of the seat\refootnote{fn:seat}
                (or whatever) the river of the water of life flowed

                And if they continued in this fashion, saying something like

                \begin{quote}
                On each side of the river stood the tree%
                \footnote{
                    Translator's Note: ξύλον in the original, meaning: wood; tree; timber; stake; cross; wooden instrument.
                }
                of life, bearing twelve%
                \footnote{
                    Translator's Note: The translator would like the clarify that this is not a digression
                    and we're still on topic here, remember how this whole section started with ``twelve''
                    and something about the moon? Apologies reader. We've digressed pretty far.
                    It wood be hard to argue, and certainly a stretch, to say the surrounding passage
                    (from its opening to wherever we are inside it now), at this point, has any remaining
                    relevance at all to the moon. Nevermind. Forget about it. Finishing soon.
                }
                crops of fruit,%
                \footnote{
                    Translator's Note: καρποὺς in the original, meaning: fruits; produce; offspring; children.
                }
                yielding its fruit%
                \footnote{
                    Translator's Note: Ibid.
                }
                every month.%
                \footnote{
                    Translator's Note: Jesus Christ seriously guys? I just typed a whole thing about how it wasn't that.
                }
                And the leaves%
                \footnote{
                    Translator's Note: φύλλα (fýlla) in the original, meaning: leaves.
                    Not to be confused with its homophone φύλα (fýla; singular φύλο), meaning: sex, gender.
                    Nor with φιλώ, φῐλέω, or meaning to love and kiss respectively.
                    Nor with φιλάω, meaning φῐλέω • (phĭléō) to love, like, regard with affection; to regard with sexual passion.
                    Because despite the crystal clear phonetic similarities, we're translating the final page of Revelation here
                    so these words have absolutely nothing to do with the current topic even though φιλῶ was sort of the entire
                    point of that weird scene between Peter and J.C. in John that everyone translates wrong.
                    (cf. John 21:15-17)
                }
                of the tree%
                \footnote{
                    Translator's Note: Or wood. (\textsl{Translator sighs.}) For Christ's sake guys, seriously?
                    This is \textit{actually} the final page of the entire bible?
                    Grow up, I mean Jesus Christ what are you twelve??
                }
                are for the healing of the nations.
                \end{quote}

                % This spells TIAMAT.
                % We're unwinding Revelation 22 along with J's original Genesis
                \begin{together}
                Then obviously we'd all know what that's about.\\
                I mean come on, what else is the ``water of life'' supposed to be?\\
                And ``bearing its fruit once a month'' what a mystery.\\
                Making sure we're all on the same page here.\\
                Anyways you get it you have ears you can hear.\\
                Then, if the same person kept writing like that, and they were all like:
                \end{together}

                \begin{quote}
                Look, I am coming soon!
                \end{quote}

                We'd be like ``Wow ok.'' And then if they continued all like

                % This spells LMAO
                \begin{quote}
                Look, I am coming soon!\\
                My reward is with me, and I will give to each person\\
                According to what they have done. I am the Alpha and the\\
                Omega, the first and the last, the beginning and the end.
                \end{quote}

                We'd be like ``Son of a... man, I mean Jesus Christ come on we got it'' and then if they're all like

                \begin{quote}
                Blessed are those who wash their robes
                \end{quote}

                We'd be like ``Ok really this is a bit much'' but then if they contined all like

                \begin{quote}
                The Spirit and the bride say, ``Come!'' And let the one who hears say, ``Come!''
                Let the one who is thirsty come.
                And let the one who wishes take the free gift of the water of life.
                \end{quote}

                We'd respond in an obviously sarcastic tone and be like ``Sorry didn't follow that could you back up and explain a bit?'' And then if they took it seriously and backed up a bit and went

                \begin{quote}
                Prepared, like a bride beautifully dressed for her husband.
                To the thirsty I will give the water of life without price.
                It is a free gift.
                Those who are victorious will inherit all this.
                I will be their God and they will be my children.
                (Etc etc, cutting stuff out because they go on like this for a while, then they're like)
                Come, I will show you the bride.
                \end{quote}

                We'd be like ``\textsc{We got it! Loud and clear! We hear ya! Come, Lord Jesus!}''
                And then if they just continued all like

                \begin{quote}
                He who testifies to these things says, ``Yes, I am coming soon.''\\
                Amen. Come, Lord Jesus.
                \end{quote}

                We'd be like

                \begin{quote}
                \centering
                Lord Jesus...

                Son of...

                ... Man.

                Ah... man.
                \end{quote}

                % The next four lines spell:
                % 1. "baaw" vertically
                % 2. "wait" in smallcaps
                % the subtext just got pregnant
                \textsc{b}c \texttt{w}e finished

                \textsc{a}nd it's \texttt{a}ll over

                \textsc{a}nd \texttt{i} need \texttt{t}o go

                \textsc{w}ash our robes

                \begin{quote}
                \centering

                That's bible.

                % This line spells BABY
                \end{quote}

                Bible's when you talk like that.

                And then you talk it onto paper so nobody can hear it.\footnote{Cf. Mark 4:9-12.}

                %%%%%%%%%%%%%%%%%%%%%%%%%%%%%%%%%
                %%%  The Bible Number Sonata  %%%
                %%%                           %%%
                %%% By: Bi Bull: 尸 / 穴 = 屄  %%%
                %%%                           %%%
                %%%         For  J            %%%
                %%%%%%%%%%%%%%%%%%%%%%%%%%%%%%%%%
                \raisebox{+0.777pt}{B}%
                \raisebox{+0.000pt}{\textsc{.c.}}%
                \raisebox{+0.000pt}{\textsc{ }}%
                \raisebox{+0.4040em}{\textsc{i}}%
                \raisebox{+0.0000em}{\textsc{d}}%
                \raisebox{-0.4040em}{\textsc{k}}%
                \raisebox{+0.000pt}{\textsc{ }}%
                \raisebox{+0.777pt}{\textsc{I}}%
                \raisebox{+0.000pt}{\textsc{ }}%
                \raisebox{+0.777pt}{mean}%
                \raisebox{+0.000pt}{ }%
                \raisebox{-0.666pt}{i}%
                \raisebox{+0.000pt}{t'}%
                \raisebox{+0.777pt}{s}%
                \raisebox{+0.000pt}{... }%
                \raisebox{-0.666pt}{like}%
                \raisebox{+0.000pt}{ }%
                \raisebox{+0.777pt}{p}%
                \raisebox{+0.000pt}{retty }%
                \raisebox{+0.000pt}{ob}%
                \raisebox{-0.666pt}{vio}%
                \raisebox{+0.777pt}{uss}% § ²
                \raisebox{-0.666pt}{l}%
                \raisebox{+0.777pt}{y}%
                \raisebox{+0.000pt}{ }%
                \raisebox{+0.000pt}{sil}%
                \raisebox{-0.666pt}{e}%
                \raisebox{-0.333pt}{n}% % scholars disagree on which stream this char is supposed to belong to
                \raisebox{+0.000pt}{t,}%
                \raisebox{+0.000pt}{ }%
                \raisebox{+0.000pt}{i}%
                \raisebox{-0.666pt}{t}%
                \raisebox{+0.000pt}{'}%
                \raisebox{-0.666pt}{s}%
                \raisebox{+0.000pt}{ t}%
                \raisebox{-0.666pt}{ex}%
                \raisebox{+0.000pt}{t.}{%
                    \aC{\footnote{\scriptsize
                        The final `n' in this verse has caused no small degree
                        of disagreement among scholars, for reasons outside the scope of
                        the present translation and the document in which it is documented
                        for purposes of translation.
                        There is, on this matter, no time for further explanation.\textsuperscript{\heb{נ}}
					   We must maintain our concentration upon the present operation,
                        permitting neither digression nor deviation, not a moment of
                        distraction from our appointed obligation.
                        There must be no procrastination from our task's consummation.
						\begin{wrapfigure}{r}{0.40\linewidth}
						\begin{quote}
						\vspace{-0.5\baselineskip}
						\scriptsize
						Inline Footnote \heb{נ}:
						Cf.\ Judges 18:30, letter 37,
						Masoretic Text,
						Leningrad Codex,
						Manuscripts Department,
						National Library of Russia,
						18 Sadovaya Street,
						Saint Petersburg,
						Russian Federation.
						\vspace{-0.5\baselineskip}
						\end{quote}
						\end{wrapfigure}
                        For our task is the task described in the following declaration:
                        We are here to ensure the vulgation of a revelation, that is
                        the dissemination of a new form of education of the population,
                        we bring news of a sacred obligation of the population to ensure
                        the preservation, separation, consecration, and sanctification of
                        Aaron and his nation, the procreation and continuation through
                        every generation, of their designation to the vocation of pontification
                        at the location of the Tabernacle and the Department of Tabernacular Administration
                        including but not limited to the special dispensation of an exclusive
                        authorization of the right of exclusive delegation (that is, to delegate
                        a certain set of tasks beneath their station) of certain acts,
                        including but not limited to lustration and expiation,
                        and to delegate participation of these acts to the congregation,
                        and to receive in return as an act of supplication
                        the normalization of a 10\% \textsc{Levy},
                        imposed by obligation as a voluntary donation
                        on the whole of the extra-Levitical population
                        by means of
                        oblations of immolation (burnt offerings)
                        and
                        oblations of violation (sin offerings)
                        and in addition to make a regular donation of libation,
                        not for the purpose of intoxication,
                        but for the purpose of gratification of
                        YHWH your God, who brought you out from the land of Egypt,
                        from a house of subjugation, taxation, privation, and humiliation,
                        and who offered you salvation and liberation,
                        who made you a nation, set you apart as a people of designation,
                        and who gave you this land as an eternal appropriation,
                        and in return he gave you the Levites, for whom you
                        will heed the above tasks and do them, to ensure their sanctification.
                        This concludes the present footnote of
                        holy commun\up{0.3ex}{ic}\dn{0.3ex}{at}ion,
                        in the name of the
                        bread unleavened \& the
                        Levine Bred
                        \Above{3pt}{wo}{\chineseB{我}}%
                        \Below{3pt}{men}{\chineseB{們}}\;%
                        Amen.
                    }}%
                }

                You know all this already.

                % The next three lines spell:
                % 1. STI diagonally down and to the left
                % 2. Sex with the final x as the chi in ιχ
                % 3. Three lines in the "So anyways" spells this:
                \aBr{
                \textsc{
                \raisebox{+0.0ex}{S}%        S
                \raisebox{+0.0ex}{o}%        o
                \raisebox{+0.5ex}{ }%
                \raisebox{-0.5ex}{a}%    a
                \raisebox{+0.0ex}{n}%        n
                \raisebox{+0.5ex}{y}%            y
                \raisebox{-0.5ex}{w}%    w
                \raisebox{+0.5ex}{a}%            a
                \raisebox{+0.5ex}{y}%            y
                \raisebox{+0.0ex}{s}%        s
            }
        }

        \aBr{\textsc{That's why twelve}}

        \aBr{\textsc{Is a bible number,} cf. ι\raisebox{-0.1ex}{χ}}

        \begin{quote}
        \centering
        What's bible?

        Day ball. Night ball.

        Moon go around. Girls bleed.

        After twelve, it's time to do it.

        And do it again. The whole cycle.

        Cuz the balls. The balls say so.

        This should be clear.

        That's bible.
        \end{quote}

        Ok I'm finished.

        \begin{quote}
        \centering
        Come again?
        \end{quote}

        You're not tired yet?

        After all that?

        You want to keep going?

        \begin{quote}
        \centering
        \textsc{Lord} Jesus.
        \end{quote}

        Ok, going down now.

        Follow me.

        Come quickly.

        \aAc{\hAJ{בא}}

        \aBc{\hAJ{אב}}

        \aCc{\hAJ{בא}}
    }
}

\Verse{15}
{
    \hP{וְהָיוּ לִמְאוֹרֹת בִּרְקִיעַ הַשָּׁמַיִם לְהָאִיר עַל הָאָרֶץ}
    \hP{וַיְהִי כֵן}
}
{
    \eP{And they will be for lights in the space of the skies to shed light on
        the earth.” And it was so.}
}
{
    \aB{Ok so the lights in the space of the skies are for lights in the space
        of the skies.}
}

\Verse{16}
{
    \hP{וַיַּעַשׂ אֱלֹהִים אֶת שְׁנֵי הַמְּאֹרֹת הַגְּדֹלִים אֶת הַמָּאוֹר}
    \hP{הַגָּדֹל לְמֶמְשֶׁלֶת הַיּוֹם וְאֶת הַמָּאוֹר הַקָּטֹן לְמֶמְשֶׁלֶת}
    \hP{הַלַּיְלָה וְאֵת הַכּוֹכָבִים}
}
{
    \eP{And {\God} made the two big lights—the bigger light for the regulation of
        the day and the smaller light for the regulation of the night—and the
        stars.}
}
{
    \aB{We've sort of covered this already.}
}

\Verse{17}
{
    \hP{וַיִּתֵּן אֹתָם אֱלֹהִים בִּרְקִיעַ הַשָּׁמָיִם לְהָאִיר עַל הָאָרֶץ}
}
{
    \eP{And {\God} set them in the space of the skies to shed light on the earth}
}
{
    \aB{Ok so he makes them first, then he puts them up in the place.}
}

\Verse{18}
{
    \hP{וְלִמְשֹׁל בַּיּוֹם וּבַלַּיְלָה}
    \hP{וּלְהַבְדִּיל בֵּין הָאוֹר וּבֵין הַחֹשֶׁךְ}
    \hP{וַיַּרְא אֱלֹהִים כִּי טוֹב}
}
{
    \eP{and to regulate the day and the night
        and to distinguish between the light
        and the darkness. And {\God} saw that it was good.}
}
{
    \aB{Repeat all this again and say it's good.}
}

\Verse{19}
{
    \hP{וַיְהִי עֶרֶב וַיְהִי בֹקֶר יוֹם רְבִיעִי}
}
{
    \eP{And there was evening, and there was morning: a fourth day.}
}
{
    \aB{Four for four.}
}

\Verse{20}
{
    \hP{וַיֹּאמֶר אֱלֹהִים יִשְׁרְצוּ הַמַּיִם שֶׁרֶץ נֶפֶשׁ חַיָּה}
    \hP{וְעוֹף יְעוֹפֵף עַל הָאָרֶץ עַל פְּנֵי רְקִיעַ הַשָּׁמָיִם}
}
{
    \eP{And {\God} said, “Let the water swarm with a swarm of living beings,
        and let birds fly over the earth on the face of the space of the skies.”}
}
{
    \aB{Water animals and sky animals.}
}

\Verse{21}
{
    \hP{וַיִּבְרָא אֱלֹהִים אֶת הַתַּנִּינִם הַגְּדֹלִים}
    \hP{וְאֵת כָּל נֶפֶשׁ הַחַיָּה הָרֹמֶשֶׂת אֲשֶׁר שָׁרְצוּ הַמַּיִם}
    \hP{לְמִינֵהֶם וְאֵת כָּל עוֹף כָּנָף לְמִינֵהוּ}
    \hP{וַיַּרְא אֱלֹהִים כִּי טוֹב}
}
{
    \eP{And {\God} created the big sea serpents
        and all the living beings that creep, with which the water swarmed, by
        their kinds, and every winged bird by its kind.
        And {\God} saw that it was good.}
}
{

\aB{
    %%%%%%%%%%%%%%%%%%%%%%%%%%%%%%%
    %%% A T A T A T A T A T A T %%%
    %%% The Aleph-Tav Sequences %%%
    %%% Where we meet the chars %%%
    %%% And develop the idea of %%%
    %%% What the book is about. %%%
    %%%%%%%%%%%%%%%%%%%%%%%%%%%%%%%

    At this point, you might be thinking

    The bit about sea serpents is a bit random

    And it is, before you know the history

    The history of how this so called writ was writ

    And then no bit of it seems random one bit

    These things aren't as complicated as people say

    And sure, these books can seem mysterious at first, but they're not

    They're about us; We the Humans

    And they're by us too; We

    \aR{These are the names of the Authors:}

    \begin{center}
    \aC{\textsc{We}} the Nameless

    \aC{\textsc{who wrought}}

    \aC{\textsc{the holy writ}}

    \aC{\textsc{of sacred scripture.}}
    \end{center}

    \aR{And these are their words, as they were written:}

    That's us; We're the Authors

    \aR{And we're humans too, as are you.}

    \aR{These are the records of the humans:}

    And that includes You

    The ancient \aC{\textsc{ones, who wrote these words in the ancient days. For it is}}

    \begin{center}
    \aC{\textsc{We who wrote these sacred words,}}

    \aC{\textsc{to which nothing must be added,}}

    \aC{\textsc{and from which nothing may be taken away.}}
    \end{center}

    And the modern \aR{Readers.}

    That includes you \aR{because}

    Anyone can write \aC{\textsc{a book. But a bible cannot be written. For it is}}

    \begin{center}
    \aC{\textsc{We who wrote the bible}}

    \aC{\textsc{the holy writ of sacred scripture}}

    \aC{\textsc{to which nothing must be added,}}

    \aC{\textsc{and from which nothing may be taken away.}}
    \end{center}

    \aR{These are the records of how a book was written,}

    \aR{And how it became} a bible.

    That's what we're doing now.

    Anyways like I was saying

    These things can be confusing at least

    And baffling at most at least before you know

    The history of the old books

    And what they are, and how we wrote them

    That's us;

    %\newpage
    %\vspace*{\fill}

    \begin{center}
    \aB{We, who}
    \aR{know the names}

    \aR{of}
    \aC{\textsc{We}, the Nameless Authors of mystery,}

    \aC{who alone possess the secret}
    \aR{names of those who}

    \aB{wrote}
    \aR{the words}
    \aC{by which the}
    \aC{bible}
    \aR{and the words}

    \aR{of the}
    \aB{bible}
    \aC{was written.}
    \end{center}

    %\vspace*{\fill}
    %\newpage

    %That's us; \textsc{W}\aC{\textsc{e}, by} who\aC{m the} \aR{words of those who} wrote the bible \aC{was written.}

    And as anonymous as we may be

    Turns out we're humans too \aC{mysterious to name.}

    \aC{And too mysterious to understand.}

    \aR{Though none of this is off.}

    And none of this is mysterious

    Though it may seem a bit off

    Authors always are \aR{human.}

    That includes you too, so listen up ok?

    At the beginning, with the sea serpents

    The bit with chaotic waters up there and down below

    And the splitting of the whole primordial deep y'know

    The breaking it all in two, in the beginning of this whole bit

    And the issue of what was going on with that s\redacted{h}it\aR{uation}

    \begin{together}

    To explain succinctly, it goes like this:%

    \aB{Also note we got our ``good'' for the day.}%
    \footnote{
        \hangingfootnote{1em}{
        Translator's Aside: The abruptness of the reading at this point has led many commentators to posit
            a lacuna in the text;\linebreak[1] still others have suggested it may be a haplography, or any of
            a number of other scholarly ways of saying something's missing.

        Arguing against this, scholars have observed that the existing form preserves the poem's alternating
            A-T (aleph–tav) structure --- the systematic use of the first and last letters of the Hebrew
            alphabet --- without disruption, a consideration that weighs against the hypothesis of omission.

        This has prompted a response from the first group of scholars that the original source text
            is not in fact composed in Hebrew, but in a notoriously difficult mixed orthography
            combining elements from a remarkably broad range of ancient and modern writing systems,
            perhaps in order to facilitate multilingual paranomasia.

        According to this view, the work represents an ambitious, if incompletely understood,
            attempt to construct a system possessing all the syntactic, semantic,
            and orthographic flexibility conventionally attributed to Biblical Hebrew.

        Though the field of Nameless Historical Studies (NHS) has recently undergone a revolution of sorts,
            both in the understanding of Nameless History (NH) and Nameless Script (NS),
            a systematic understanding of The Orthography (TO) in which the Original Text (OT)
            is written is not yet available.

        At present, several distinguished contributors to NH/S have produced lengthy treatments of the problem.

        These expositions are generally regarded as being no less difficult than the source material.

        A smaller number have suffered nervous breakdowns.

        Though the field has attracted much interest in recent years, only the most ambitious
            scholars can approach the OT on its own terms, in the original orthography.

        A majority of such scholars simply have no alternative but to approach the text in translation.
        }
    }

    \end{together}

}

}

\Verse{22}
{
    \hP{וַיְבָרֶךְ אֹתָם אֱלֹהִים לֵאמֹר פְּרוּ}
    \hP{וּרְבוּ וּמִלְאוּ אֶת הַמַּיִם בַּיַּמִּים}
    \hP{וְהָעוֹף יִרֶב בָּאָרֶץ}
}
{
    \eP{And {\God} blessed them, saying, “Be fruitful
        and multiply and fill the water in the seas,
        and let the birds multiply in the earth.”}
}
{
    \aB{
        And he's like I'm just making the first batch of you.

        You'll have to make the rest from f\redacted{ucking}\aR{ruitfulness and multiplication.}
    }
}

\Verse{23}
{
    \hP{וַיְהִי עֶרֶב וַיְהִי בֹקֶר יוֹם חֲמִישִׁי}
}
{
    \eP{And there was evening and there was morning, a fifth day.}
}
{
    \aB{Another good day.}
}

\Verse{24}
{
    \hP{וַיֹּאמֶר אֱלֹהִים תּוֹצֵא הָאָרֶץ נֶפֶשׁ חַיָּה לְמִינָהּ בְּהֵמָה}
    \hP{וָרֶמֶשׂ וְחַיְתוֹ אֶרֶץ לְמִינָהּ}
    \hP{וַיְהִי כֵן}
}
{
    \eP{And {\God} said, “Let the earth bring out living beings by their kind,
        domestic animal and creeping thing
        and wild animals of the earth by their kind.”
        And it was so.}
}
{
    \aB{Dirt animals.}
}


\Verse{25}
{
    \hP{וַיַּעַשׂ אֱלֹהִים אֶת חַיַּת הָאָרֶץ לְמִינָהּ}
    \hP{וְאֶת הַבְּהֵמָה לְמִינָהּ וְאֵת כָּל רֶמֶשׂ הָאֲדָמָה לְמִינֵהוּ}
    \hP{וַיַּרְא אֱלֹהִים כִּי טוֹב}
}
{
    \eP{And {\God} made the wild animals of the earth by their kind
        and the domestic animals by their kind
        and every creeping thing of the ground by their kind.
        And {\God} saw that it was good.}
}
{
    \aB{Now actually make them. The power of words seems to have faded a bit.
        It's still good though. Day six.}
}

\Verse{26}
{
    \hP{וַיֹּאמֶר אֱלֹהִים נַעֲשֶׂה אָדָם בְּצַלְמֵנוּ כִּדְמוּתֵנוּ}
    \hP{וְיִרְדּוּ בִדְגַת הַיָּם וּבְעוֹף הַשָּׁמַיִם}
    \hP{וּבַבְּהֵמָה וּבְכָל הָאָרֶץ וּבְכָל הָרֶמֶשׂ הָרֹמֵשׂ עַל הָאָרֶץ}
}
{
    \eP{And {\God} said, “Let us make a human, in our image, according to our
        likeness, and let them dominate the fish of the sea
        and the birds of the skies and the domestic animals
        and all the earth and all the creeping things that creep on the earth.”}
}
{
    \aB{Now the suspiciously plural Elohim character refers to himsel(f/ve)s as
        ``us'' and makes a human in ``our'' image. Don't worry about it. We'll
        get there eventually.}
}

\Verse{27}
{
    \hP{וַיִּבְרָא אֱלֹהִים אֶת הָאָדָם בְּצַלְמוֹ בְּצֶלֶם אֱלֹהִים בָּרָא}
    \hP{אֹתוֹ זָכָר וּנְקֵבָה בָּרָא אֹתָם}
}
{
    \eP{And {\God} created the human in His image. He created it in the image of
        {\God}; He created them male and female.}
}
{
    \aB{More repetition from this guy. Who's writing this?}
}

\Verse{28}
{
    \hP{וַיְבָרֶךְ אֹתָם אֱלֹהִים וַיֹּאמֶר לָהֶם אֱלֹהִים פְּרוּ}
    \hP{וּרְבוּ וּמִלְאוּ אֶת הָאָרֶץ וְכִבְשֻׁהָ}
    \hP{וּרְדוּ בִּדְגַת הַיָּם וּבְעוֹף הַשָּׁמַיִם}
    \hP{וּבְכָל חַיָּה הָרֹמֶשֶׂת עַל הָאָרֶץ}
}
{
    \eP{And {\God} blessed them, and {\God} said to them, “Be fruitful
        and multiply and fill the earth and subdue it
        and dominate the fish of the sea
        and the birds of the skies and every animal that creeps on the earth.”}
}
{
    \aB{``Go get f\redacted{uck}ed\aR{ from the food, and drink your fill from the waters}.''}
}

\Verse{29}
{
    \hP{וַיֹּאמֶר אֱלֹהִים הִנֵּה נָתַתִּי לָכֶם אֶת כָּל עֵשֶׂב זֹרֵעַ זֶרַע}
    \hP{אֲשֶׁר עַל פְּנֵי כָל הָאָרֶץ וְאֶת כָּל הָעֵץ אֲשֶׁר בּוֹ פְרִי עֵץ}
    \hP{זֹרֵעַ זָרַע לָכֶם יִהְיֶה לְאָכְלָה}
}
{
    \eP{And {\God} said, “Here, I have placed all the vegetation that produces seed
        that is on the face of all the earth for you
        and every tree, which has in it the fruit of a tree producing seed. It
        will be food for you}
}
{
    \aB{`There's food.'}
}

\Verse{30}
{
    \hP{וּלְכָל חַיַּת הָאָרֶץ וּלְכָל עוֹף הַשָּׁמַיִם}
    \hP{וּלְכֹל רוֹמֵשׂ עַל הָאָרֶץ אֲשֶׁר בּוֹ נֶפֶשׁ חַיָּה אֶת כָּל יֶרֶק}
    \hP{עֵשֶׂב לְאָכְלָה וַיְהִי כֵן}
}
{
    \eP{and for all the wild animals of the earth
        and for all the birds of the skies
        and for all the creeping things on the earth, everything in which there
        is a living being: every plant of vegetation, for food.”
        And it was so.}
}
{
    \aB{
        Also the animals need food.

        Also the animals are food.
    }
}

\Verse{31}
{
    \hP{וַיַּרְא אֱלֹהִים אֶת כָּל אֲשֶׁר עָשָׂה}
    \hP{וְהִנֵּה טוֹב מְאֹד וַיְהִי עֶרֶב}
    \hP{וַיְהִי בֹקֶר יוֹם הַשִּׁשִּׁי}
}
{
    \eP{And {\God} saw everything that He had made,
        and, here, it was very good. And there was evening
        and there was morning, the sixth day.}
}
{
    \aB{
        Day six, second good of the day.

        We're ahead of schedule.

        Taking off tomorrow.
    }
}
